\section{Why an analyser}
\label{sec:whyanalyzer}

Until now every observation has been made by one of the two halves. Your driver says that it sent
a frame; the hardware group's testbench says that the controller transmitted one. Both can be
wrong in the same way, since both are written against the same document by two people who may have
misread it identically.

A \textbf{bus analyser} is a third party that believes neither of you. It is a real CAN controller
on the bus, and it decodes what is actually on the wires.

That is why the lab comes last and not first: only now is there traffic to look at, and only now
does the answer mean anything.

\begin{aside}
\note The analyser also acknowledges your frames. That is worth knowing: a lone node on the bus
never gets an acknowledgement, and \emph{a complete} CAN controller would then have retransmitted
for all eternity. The design in this course never reads the ACK slot back, so it notices nothing
--- but the difference becomes visible in the trace as soon as the analyser is connected.
\end{aside}

\section{Reading a trace}
\label{sec:trace}

Set the analyser to the same bit rate as the design, and check that the termination is still
120~\textohm{} at each end --- not three resistors just because there are now three nodes.

Send a frame with a known ID and a known payload from your node, and find it:

\begin{booktable}{@{}lL@{}}
\thead{Column} & \thead{What you check} \\ \midrule
Identifier & Does it match what you wrote to \code{TX_ID}? \\
DLC & Does it match \code{TX_DLC}? \\
Data & Are the bytes right, \emph{and in the right order}? This is where a swap of
  \code{TX_DATA_LO}/\code{HI} finally becomes visible. \\
Timestamp & How often do they arrive? Does that match how often your code calls \code{send()}? \\
\end{booktable}

Then the other way round: send a frame \emph{from} the analyser to your node. Does
\code{receive()} return \code{true}? Is the content right? And --- the thing that is easy to
forget --- did \code{STATUS} bit 1 go low when you wrote \code{RX_ACK}?

\section{Looking for the simplifications}
\label{sec:simplifications}

The design is simplified compared with real CAN in several respects. Some of them can be seen from
the outside, and that is the most interesting part of the lab.

\begin{booktable}{@{}lL@{}}
\thead{Simplification} & \thead{What you can look for} \\ \midrule
No ACK read-back & Disconnect the analyser and transmit with only your node on the bus. What does
  the error flag say? What \emph{would} a real node have done? \\
No error frames & What happens if the analyser sends one? \\
No automatic retransmission & A frame that went wrong is not resent by the hardware. Can you see
  that? \\
No receive queue & Let the analyser send two frames close together. Do both reach your
  application? \\
\end{booktable}

\textbf{Finding a discrepancy is not the same as finding a fault.} The task is to be able to say
\emph{which documented simplification} explains what you see --- and which discrepancies have
\emph{no} such explanation. It is the latter that are interesting, for both classes.

\section{Bus load}
\label{sec:busload}

The analyser shows the bus load. Compare it with the time budget in \kapref{ch:spi}:

\begin{itemize}
  \item What is the load with your node as the only transmitter?
  \item A \code{send()} takes six SPI transactions, that is, around 270~µs. A maximal CAN frame is
    around 130~µs. \textbf{What is actually limiting you --- the bus or the link?}
\end{itemize}

The answer to that last question is the whole point of measuring instead of assuming, and it is
the one place in the course where you can see it with your own eyes.

\section{One frame, all the way}
\label{sec:endtoend}

Finish by tracing a single frame through the whole stack, naming each layer and who built it:

\begin{console}
app::EchoNode            -> you
driver::can::Interface   -> you
driver::can::Spi         -> you
driver::can::ByteTransport -> you
AvrSpiTransport          -> you
===== four wires =====
spi_slave                -> supplied
spi_reg_bridge           -> the hardware group
register_bank            -> the hardware group
can_controller           -> the hardware group
===== transceiver =====
the CAN bus              -> physics
\end{console}

Ten layers, two classes, one contract --- and a frame that arrived.
